\chapter{系统需求分析}
\section{业务场景与角色}
系统包含学生、实验室管理员和系统维护人员三类角色。学生负责检索设备、提交申请和确认归还；管理员负责审批申请、发放设备、登记归还及处理异常；维护人员负责账号、角色和基础字典配置。不同角色只能访问与职责相关的页面和接口。
\section{功能需求}
系统功能划分如下：
\begin{enumerate}
\item 账号登录与角色授权：验证身份并根据角色展示可用功能。
\item 设备档案管理：维护编号、名称、分类、存放位置、可借数量和当前状态。
\item 借用申请管理：学生填写用途、借用时间和归还时间，系统校验申请数据。
\item 审批与归还管理：管理员批准或驳回申请，领取和归还操作改变业务状态。
\item 查询与审计：按设备、申请人、状态和时间范围检索记录，并查看关键操作日志。
\end{enumerate}
\section{非功能需求}
系统应具备清晰的错误提示和可追溯的操作记录。库存更新需要在一次事务内完成，避免并发申请导致可借数量出现负值。常用列表查询应在合理数据规模下快速返回，接口错误应使用统一结构，便于前端处理。
\section{可行性分析}
系统采用浏览器访问方式，不要求学生安装专用客户端；关系数据库能够表达设备、申请和日志之间的关联；状态机和事务是成熟的软件工程方法。因此，项目在技术、运行和维护方面均具备可行性。
